\SourceLeaf{72}%
ihm hier eigenthümlichen Abänderungen zugleich eine gedrängte Übersicht der
sämmtlichen übrigen Operationen verbinde, die zur ersten Bestimmung einer Cometenbahn
erforderlich sind. Wir bemerken dabei, dass auch einige der andern
Formeln hier in einer von der sonst gewöhnlichen verschiedenen Gestalt erscheinen,
bei welcher man aber nur in der Voraussetzung noch etwas an Kürze
gewinnt, dass man sich dabei der von dem Verfasser vor einem Jahre bekannt
gemachten Tafel zur unmittelbaren Berechnung der Logarithmen von Summen
und Differenzen bedient. Man hat auf diese Weise hier auf zwei oder drei Seiten
alles beisammen, was ausser den Tafeln zu der Berechnung einer Cometenbahn
nöthig ist, und alles ist zugleich durch ein Beispiel an dem letzten Cometen
erläutert, wo der Verfasser aus seinen Beobachtungen vom 7., 14. und
21. April Elemente ableitet, welche von den oben mitgetheilten und auf die gesammten
Beobachtungen gegründeten nur unbedeutend verschieden sind.
\par\bigskip\centerline{\rule{.24\textwidth}{.3pt}}
